\section{Implementación}

Las dos implementaciones son clases genéricas de C++20 obtenidas del
repositorio del curso y conservadas en \texttt{include/range\_queries/}.
Reciben un vector de elementos de tipo \texttt{T} y consultan intervalos
semiabiertos \([i,j)\). Para obtener el índice del mínimo se utilizaron pares
\((\text{valor},\text{índice})\), cuyo orden lexicográfico resuelve empates
por el menor índice. Se reemplazaron los comentarios y se añadió la cabecera
\texttt{<bit>} a Sparse Table, sin modificar la lógica de los algoritmos.

\subsection{Sparse Table}

La clase \texttt{sparse\_table} guarda una matriz de vectores \texttt{ST}
y una función de combinación \texttt{f}. El constructor obtiene
\(m=\lfloor\log_2 n\rfloor\) mediante \texttt{\_lg}, reserva \(m+1\)
filas de \(n\) elementos y copia la entrada en \texttt{ST[0]}. Dos ciclos
construyen las filas siguientes: \texttt{p} representa la longitud de un
bloque y \texttt{q} el doble; cada posición válida combina dos bloques
consecutivos de la fila anterior.

Para RMQ se empleó \texttt{query\_idempotent}, que calcula el nivel
\texttt{b} a partir de la longitud del intervalo y combina los bloques
de esa fila que comienzan en \(i\) y \(j-2^b\). La clase también
incluye \texttt{query}, que recorre bloques disjuntos, pero no es la variante
seleccionada para esta comparación. Como no dispone de \texttt{update},
en el escenario dinámico se reconstruyó desde el arreglo modificado.
La implementación proviene del repositorio del curso \cite{edaa-sparse-table}.

\subsection{Segment Tree}

La clase \texttt{segment\_tree} almacena la estructura en un vector
\texttt{ST} de \(2n\) elementos. El constructor copia la entrada en las
hojas, desde \texttt{ST[n]} hasta \texttt{ST[2n - 1]}, y recorre los
índices internos en orden descendente para guardar el mínimo de cada par
de hijos mediante \texttt{\_min}.

El método \texttt{query} desplaza los extremos del intervalo a las hojas e
inicializa \texttt{q} con un elemento del rango. Un ciclo asciende por el
árbol, incorpora los nodos de los bordes según la paridad de los extremos
y conserva el menor candidato. El método \texttt{update} cambia una hoja
y recalcula el mínimo de sus ancestros, combinando cada nodo con su hermano
hasta alcanzar la raíz, sin reconstruir el vector.
La implementación proviene del repositorio del curso \cite{edaa-segment-tree}.
